\chapter{The accounting language}
\label{ch:language}

\section{What the current language can express}

A company model needs to keep disclosures, possible histories and chosen
methods distinct as its accounts evolve. A reported balance constrains a
history; it does not reveal every contract or transaction that generated it.
A collection policy supplies a behavioural assumption; it does not establish
what a customer will pay. An accounting method turns an eligible event into
state changes and postings; a balanced result alone does not establish that
the method applies. The language makes these distinctions explicit so that
reports retain the meaning of the facts and assumptions behind them.

The current design has several checked, bounded language profiles. Together
they cover source-constrained histories, conditional company paths, selected
shared-book operations and complete method authoring. Their interfaces are
explicit, but they are not yet one complete whole-company grammar.
Section~\ref{sec:common-program-authoring} shows the present composition
boundary: a single program contains complete method text and registries,
selects an inspectable state-invariant definition, and retains full event
intent beside the inputs used in a specialized calculation. Invariant
algorithms are still imported typed rules, and a selected native event
contract retains its own implementation and applicability limits.

To represent the historical record, company method selection first binds to
a source-constrained partial state
(Section~\ref{sec:company-source-state}). Statement observations have
distinct relation roles rather than one undifferentiated set of totals
(Section~\ref{sec:statement-context-relations}). Selected balances are tested
jointly under explicit rounding and identity hypotheses
(Section~\ref{sec:joint-source-measurements}), with opening, movement and
closing dates kept separate across statements
(Section~\ref{sec:temporal-statement-network}). These declarations can expose
inconsistency or missing information without inventing the underlying books.

Several histories can remain consistent with those observations.
Section~\ref{sec:conditional-history-families} represents finite conditional
families with symbolic openings and explicit measurement attachment.
Typed economic quantities and constraints at each event give those histories
economic content (Section~\ref{sec:typed-economic-history}); cash and
quantity observations attach to the same history
(Section~\ref{sec:typed-history-observation-attachment}).
Unknown monetary amounts remain linked to selected method prerequisites
without choosing a single witness
(Section~\ref{sec:conditional-method-families}).
The retained return example follows every branch through settlement and sale,
including funding and intermediate observations
(Section~\ref{sec:native-conditional-sequence}). Outcome queries preserve
which results are possible or necessary and their exact finite bounds
(Section~\ref{sec:family-outcome-queries}). Stored definitions and results
can be recovered across restart, while the archive distinguishes stored
evidence from verified execution
(Section~\ref{sec:simulator-evidence-archive}).

Explicit future assumptions then determine how a conditional company path
evolves. The earliest matched text frontend handles the declared debt fragment.
The funded-year extension adds identified claim allocations, credit terms
and dated events; its imported annual projection still supplies the other
operating and accounting rules (Section~\ref{sec:funded-year-authoring}).
The operating-input example instead derives purchases, sales, claims,
expenses and closing statements from a historical opening, authored
quantities, prices, asset cohorts and contract terms
(Section~\ref{sec:company-input-authoring}). It imports no forecast journals:
a price or service-date change has consequences in the calculation under the
example's stated method restrictions.

Identified operations can repeat and carry unpaid claims and asset cohorts
into later years. Northstar's three-period path includes revised assumptions
and an interrupted year (Section~\ref{sec:company-operation-paths}).
Its text frontend preserves the same typed records, named settlement
allocations and incomplete paths
(Section~\ref{sec:company-path-authoring}).
Named, unweighted scenarios keep the alternative assumptions separate.
Their collection policies have pure typed expressions and bounded reusable
functions with matched text declarations
(Section~\ref{sec:company-family-language}). Function calls preserve units,
company context and the distinction between declared and actually used
inputs. They reuse a calculation without supplying missing facts or creating
endogenous events.

Combining methods requires compatible economic state as well as account names.
Portable module identities expose the additional state required by a
shared-book return method (Section~\ref{sec:company-module-language});
its matched authoring and finite execution appear in
Section~\ref{sec:company-return-reference}. Heterogeneous lease rights require
explicit account allocations and funding conditions
(Section~\ref{sec:lease-shared-book}). Aggregate balances may leave those
allocations unidentified, producing different complete or interrupted
outcomes (Section~\ref{sec:lease-allocation-family}).
Actual rounded lease disclosures constrain some possibilities before
individual contract state is known
(Section~\ref{sec:henkel-lease-measurements}).

A common caller does not establish a common book. The return-and-lease example
first establishes compatible state, then follows both kinds of event through
one declared cash balance (Section~\ref{sec:shared-return-lease-book}).
Its recorded statements reconcile while retaining unresolved tax assessment
and incomplete fiscal coverage
(Section~\ref{sec:recorded-shared-statements}). The complete opening and
mixed event sequence have matched text and typed forms
(Section~\ref{sec:shared-book-authoring}).
Section~\ref{sec:shared-tax-information} explains why the original tax
method refuses that mixed state, then develops explicitly conditional
successors for assessment, payments, period carry and annual programs.
Those successors preserve unknown declarations and stopped prefixes. They
do not establish the company's actual tax policy or complete fiscal history.

Methods state their prerequisites before applying financial effects.
Reusable selected requirements have matched typed and text declarations
(Section~\ref{sec:method-requirement-language}).
Exact input binding, dimensional products and guarded applicability require
explicit version selection (Section~\ref{sec:requirement-successors}).
Collection coverage, identity order, signed totals and state-domain findings
add constraints that scalar types cannot express
(Section~\ref{sec:collection-state-requirements}).
Typed calculations and atomic state updates then support exact source-bound
replay of the retained return example
(Section~\ref{sec:typed-return-transitions}).
Complete method text includes these prerequisites, monetary calculations,
postings and memory changes
(Section~\ref{sec:complete-method-authoring}).
The common program preserves that source and its explicit owners; it does
not silently replace a newly written method with a named demonstration.

The historical AccDSL Core 0 proposal remains the teaching baseline. It
describes entities and books, accounts, observations, shared unknowns,
constraints, alternative interpretations, contract instances, accounting
rules and reporting mappings, with pure expressions and a restricted posting
effect. Appendix~\ref{app:contracts} preserves its grammar and examples.
The proposed portable records and first company profile refine that baseline.
Numerical and forecasting algorithms can live in typed plugins; unrestricted
numerical code need not be reimplemented in the DSL. General recursive or
simultaneous algorithms, endogenous decisions, source admission and broader
accounting profiles remain separate language and method work.

The historical Core 0 compiler has six ordered stages: parsing with source
locations; name and module resolution; type and temporal-context checking;
lowering to the IR; constraint and rule-applicability analysis; and
execution-plan construction. It never silently supplies an economic
assumption to make a program compile. A valid historical representation can
still require additional assumptions before a particular simulation can run.
The proposed portable compilation sequence is specified separately in
Section~\ref{sec:portable-authoring}. Structural compilation of a current
program likewise does not execute its financial methods or resolve unknown
eligibility findings.

The IR and evaluator define meaning separately from surface spelling.
Formatting that preserves the lowered method preserves its semantic digest.
The complete authored program also retains the original method source, so its
program digest changes when those bytes change. In the checked common-program
example, adding one space leaves the lowered method and materialized financial
request identical while changing the source and program digests. Source spans
and reviewer decisions can change evidence provenance without changing the
numerical model. Evidence, semantic-model and execution identities therefore
remain distinct; a field named \code{hash} must identify what it actually binds.

For the preserved Core 0 examples, the complete grammar is
\code{spec/core-0.lark}; the original design directory retains
\code{spec/core-ir.schema.json} and \code{spec/core-contract.md}.
They define reference resolution, expression arity, type and context checks,
operations, failures and identity. A structural schema check alone does not
establish those semantic properties. This revision uses ASCII identifiers and
UTF-8 strings, with explicit lexical rules for currency suffixes, dates and
decimal literals. Parsing the complete demonstration checks the grammar; it
does not establish an implemented compiler or solver. The following type
distinctions apply before any of those operations can be interpreted.

